\section{Dask Auto-Chunking}
\label{sec:idp-dask-auto-chunking}

Dask~\cite{rocklin2015dask} is a distributed computing library designed to partition large arrays into smaller blocks, known as chunks, facilitating their parallel execution across multiple workers.
At the heart of Dask’s data partitioning strategy is the \emph{auto-chunking} feature, which attempts to automatically determine suitable chunk sizes based on heuristic-driven approaches.
These heuristics typically consider several factors:

\begin{itemize}
    \item The total size of the array in bytes.
    \item A user-configurable \texttt{target\_size} parameter (commonly set around $128\text{MB}$ per chunk).
    \item Approximate considerations regarding network bandwidth and disk \ac{I/O} operations.
\end{itemize}

The auto-chunking mechanism evaluates the dimensions of an array and partitions it into subarrays with sizes that closely match the specified target memory threshold (see Algorithm~\ref{alg:dask-autochunking}).

\begin{algorithm}[htbp]
    \caption{Dask auto‐chunking heuristic: computes the total size of the array in bytes, determines the number of chunks needed based on the target memory threshold, and partitions the first dimension accordingly while retaining full size for the other dimensions to form the chunk shape.}
    \label{alg:dask-autochunking}
    \begin{algorithmic}[1]
        \REQUIRE Array dimensions $\textit{dims} = [d_1, d_2, \dots, d_n]$, element size $s$ (bytes), target chunk size $T$
        \ENSURE Chunk shape $\textit{chunk\_shape}$
        \STATE $\textit{total\_elements} \gets \prod_{i=1}^{n} d_i$
        \STATE $\textit{total\_bytes} \gets \textit{total\_elements} \times s$
        \STATE $\textit{num\_chunks} \gets \left\lceil \frac{\textit{total\_bytes}}{T} \right\rceil$
        \STATE $\textit{size\_per\_chunk} \gets \left\lceil \frac{d_1}{\textit{num\_chunks}} \right\rceil$
        \FOR{$i \leftarrow 1$ to $n$}
        \IF{$i = 1$}
        \STATE $\textit{chunk\_shape}[i] \gets \min(d_i,\, \textit{size\_per\_chunk})$
        \ELSE
        \STATE $\textit{chunk\_shape}[i] \gets d_i$
        \ENDIF
        \ENDFOR
        \RETURN $\textit{chunk\_shape}$
    \end{algorithmic}
\end{algorithm}

In simpler, one-dimensional scenarios, this heuristic approach frequently yields effective and efficient chunk sizes.
However, the heuristic strategy often struggles with higher-dimensional datasets common in complex computations such as seismic analysis and tensor-based operations.

In multi-dimensional contexts, the simplistic heuristic frequently results in uneven or excessively small chunks.
An overly fragmented array leads to a proliferation of small tasks that dramatically increase overhead, placing undue strain on Dask’s task scheduler.
This excessive fragmentation can cause task contention and scheduling bottlenecks, negatively impacting the overall computational throughput and efficiency.

Furthermore, Dask’s default chunking heuristic operates without explicit consideration of algorithm-specific complexities or intermediate data expansions typical in sophisticated computational workflows.
Tensor-based operations, such as gradient computations, multi-stage transformations, or operator pipelines, often introduce intermediate data structures that significantly increase memory requirements beyond the original data size.
For instance, processing a moderate-sized volume like $600\times600\times600$ could inadvertently result in either numerous inadequately small chunks or excessively large chunks, the latter potentially causing \ac{OOM} errors in memory-constrained environments.

This fundamental limitation highlights the necessity for more sophisticated chunk-sizing strategies that account explicitly for the characteristics of the underlying computational tasks.
Different computational workloads impose distinct requirements; for example, data pipelines involving deep convolutional operations require carefully shaped chunks that minimize boundary overheads and control memory expansion during intermediate steps.
Conversely, simpler, volume-based operators might benefit from fewer, larger chunk divisions.

Dask’s auto-chunking mechanism currently lacks the ability to incorporate such task-specific and operator-specific information, leading to suboptimal performance in demanding scenarios typical of seismic and high-dimensional imaging applications.
Addressing this shortcoming, the subsequent section introduces \emph{memory-aware chunking}, a predictive approach that leverages an algorithm-specific memory usage model to intelligently guide chunk size selection, thereby enhancing both memory efficiency and computational performance.